\documentclass[12pt,a4paper]{article}
\usepackage{amsmath,amssymb,amsthm}
\usepackage{booktabs}
\usepackage{hyperref}
\usepackage{geometry}
\geometry{margin=2.5cm}

\newtheorem{theorem}{Theorem}[section]
\newtheorem{lemma}[theorem]{Lemma}
\newtheorem{corollary}[theorem]{Corollary}
\newtheorem{proposition}[theorem]{Proposition}
\theoremstyle{definition}
\newtheorem{definition}[theorem]{Definition}
\theoremstyle{remark}
\newtheorem{remark}[theorem]{Remark}

\title{The Quantum-to-Classical Transition in Recognition Science:\\
$J$-Cost Scaling and Einselection}

\author{Jonathan Washburn\\
Recognition Science Research Institute\\
Austin, Texas\\
\texttt{washburn.jonathan@gmail.com}}

\date{March 2026}

\begin{document}
\maketitle

\begin{abstract}
We derive the quantum-to-classical transition from Recognition
Science.  Classical behavior emerges when the $J$-cost of maintaining
quantum coherence exceeds the cost of decoherence.  For a system of
$N$ interacting ledger entries, the coherent-state $J$-cost scales
quadratically: $J_{\mathrm{ent}} = N\,j + \alpha\,N(N-1)/2$, where
$j = \varphi^{-5}$ is the single-entry cost and $\alpha > 0$ is the
pairwise entanglement cost (each of the $\binom{N}{2}$ pairs
contributing $\alpha$ to the $J$-cost).  Entangled states always
have higher cost than product states for $N > 1$.  The decoherence
timescale scales as $\tau_{\mathrm{dec}} = \tau_0 \cdot \varphi^{-Ng}$,
where $g > 0$ is the environmental coupling strength.  For
$N \sim 10^3$ and $g \sim 0.1$, this gives
$\tau_{\mathrm{dec}} \sim 10^{-65}$~s; for macroscopic objects
($N \sim 10^{23}$), $\tau_{\mathrm{dec}}$ has an exponent of
order $10^{21}$---effectively zero.  The quantum-to-classical
crossover for a 1-second observation window occurs at
$N_{\mathrm{cross}} \approx 210/g$ environmental modes.
``Pointer states'' (position, momentum) are selected by $J$-cost
minimization under environmental coupling (einselection).
Classical mechanics is the $N \to \infty$ limit.
\end{abstract}

\section{Introduction}

Why does the macroscopic world appear classical when the underlying
physics is quantum?  The decoherence program~\cite{Zurek} explains
how environmental interactions suppress superpositions, but does not
derive the decoherence timescale from first principles.
RS~\cite{RS_Axioms} does.

\section{Recognition Science Framework}
\label{sec:framework}

Recognition Science derives all physics from a single primitive,
the \emph{recognition cost functional} $J(x)$, uniquely determined
by three axioms.

\begin{definition}[RS Cost Functional]
For $x > 0$: $J(x) = \tfrac{1}{2}(x + x^{-1}) - 1$.
\end{definition}

\noindent\textbf{Axiom A1 (Normalization):} $J(1) = 0$.\\
\textbf{Axiom A2 (Recognition Composition Law, RCL):}
$J(xy) + J(x/y) = 2J(x)J(y) + 2J(x) + 2J(y)$ for all $x, y > 0$.\\
\textbf{Axiom A3 (Calibration):} $J''_{\log}(0) = 1$, where
$J_{\log}(t) := J(e^t)$.

\begin{theorem}[Cost Uniqueness]
\label{thm:unique}
The unique continuous solution to \textup{(A1)--(A3)} is
$J(x) = \tfrac{1}{2}(x + x^{-1}) - 1$.
\end{theorem}

\begin{proof}[Proof sketch]
Setting $f(t) = J_{\log}(t) + 1$ and rewriting \textup{(A2)} gives the
d'Alembert equation $f(s+t) + f(s-t) = 2f(s)f(t)$.  By the Acz\'el
classification theorem~\cite{Aczel1966}, the unique continuous solution
with $f(0) = 1$ and $f''(0) = 1$ is $f(t) = \cosh t$.
\end{proof}

\noindent The key property for the quantum-classical transition is that $J$-cost
is strictly convex ($J'' > 0$) and achieves its minimum at $x = 1$.
Nature minimizes $J$-cost: macroscopic superpositions (with quadratically
growing $J$-cost) are driven toward classical product states.  The 8-tick
epoch ($2^D = 8$ for $D = 3$) provides the fundamental time unit $\tau_0$.

\section{$J$-Cost Scaling}

\begin{definition}[$J$-Cost of $N$ Entries]
Product state: $J_{\mathrm{prod}} = N\,j_{\mathrm{single}}$
(linear in $N$).

Entangled state:
$J_{\mathrm{ent}} = N\,j_{\mathrm{single}}
+ \alpha\,N(N-1)/2$ (quadratic in $N$).
\end{definition}

\begin{lemma}[Entangled States Cost More]
For $N > 1$ and $\alpha > 0$:
$J_{\mathrm{ent}} > J_{\mathrm{prod}}$.
\end{lemma}

\begin{proof}
$J_{\mathrm{ent}} - J_{\mathrm{prod}} = \alpha\,N(N-1)/2 > 0$
for $N > 1$, $\alpha > 0$.
\end{proof}

\begin{corollary}[Cost Difference Scales Quadratically]
$\Delta J = \alpha\,N(N-1)/2 \propto N^2$.
\end{corollary}

Since nature minimizes $J$-cost, macroscopic superpositions (high
$N$, high $J_{\mathrm{ent}}$) are strongly disfavored.

\section{Decoherence Timescale}

\begin{definition}[Decoherence Time]
\begin{equation}
  \tau_{\mathrm{dec}} = \tau_0 \cdot \varphi^{-N \cdot g},
\end{equation}
where $N$ is the number of environmental modes coupled to the system
and $g > 0$ is the dimensionless coupling strength (how many $\varphi$-rungs
are shifted per environmental mode per tick).
\end{definition}

\begin{remark}[Status of $g$]
\label{rem:g-status}
The coupling $g$ is a \emph{phenomenological parameter} describing
the interaction strength between the system and environment.  It is
not yet derived from the RS Hamiltonian; a derivation of $g$ in terms
of $\varphi$ from the ledger's double-entry structure is an identified
open problem.  For a system interacting via the electromagnetic force
at room temperature, $g \sim 10^{-4}$--$10^{-2}$ provides the
observed mesoscopic decoherence timescales; for a macroscopic object
in air, $g \sim 10^{-2}$--$10^{-1}$.  When $g$ is determined from
the RS Hamiltonian, the formula $\tau_{\mathrm{dec}} = \tau_0 \varphi^{-Ng}$
will yield a parameter-free decoherence timescale.
\end{remark}

\begin{remark}[Decoherence Monotonicity]
$\tau_{\mathrm{dec}} = \tau_0 \varphi^{-Ng}$ is a decreasing
function of both $N$ and $g$: since $\varphi > 1$, increasing $N$ or
$g$ increases $Ng$, which increases $\varphi^{Ng}$ and thus decreases
$\varphi^{-Ng}$.
\end{remark}

\begin{remark}[Macroscopic Decoherence --- Order-of-Magnitude Estimate]
For macroscopic objects ($N \sim 10^{23}$, $g \sim 0.1$):
$\tau_{\mathrm{dec}} = \tau_0 \cdot \varphi^{-10^{22}}$.
Since $\log_{10}(\varphi^{10^{22}}) = 10^{22} \times 0.209
\approx 2 \times 10^{21}$, the decoherence time is
$\tau_{\mathrm{dec}} \sim 10^{-(44 + 2 \times 10^{21})}$~s---a
number with over $10^{21}$ zeros in the exponent, effectively zero.
Even for modest $N \sim 10^3$ and $g \sim 0.1$
($Ng = 100$): $\tau_{\mathrm{dec}} = \tau_0 \cdot \varphi^{-100}
\approx 10^{-44} \times 10^{-20.9} \approx 10^{-65}$~s, already
far below any observable timescale.
\end{remark}

\begin{corollary}[No Schr\"odinger's Cat]
Macroscopic superpositions decohere effectively
instantaneously.  The cat is always either alive or dead,
never both.
\end{corollary}

\section{Pointer States and Einselection}

\begin{definition}[Pointer State]
A state that is stable under environmental interaction and minimizes
$J$-cost.
\end{definition}

\begin{proposition}[Einselection from $J$-Cost --- Structural Argument]
The environment selects pointer states (position eigenstates for
massive particles, energy eigenstates for isolated systems) as the
$J$-cost minima under the environmental coupling.
\end{proposition}

\begin{remark}
The RS identification is: among all possible bases, pointer states are
those that minimize total $J$-cost when the environment is included.
Position eigenstates are pointer states because they minimize the
$J$-cost of localization: delocalized superpositions have higher
pairwise-correlation costs.  This is the RS restatement of Zurek's
einselection~\cite{Zurek}; the RS-specific content is the
cost-functional underpinning of the selection criterion.  A
derivation of pointer states as $J$-cost minima from the RS
environment--system interaction Hamiltonian is an identified open
problem.
\end{remark}

\section{Classical Limit}

\begin{proposition}[Classical Limit --- Structural Argument]
In the limit $N \to \infty$:
\begin{enumerate}
  \item $\tau_{\mathrm{dec}} = \tau_0 \varphi^{-Ng} \to 0$
  (instant decoherence; exact from the definition for any $g > 0$).
  \item Only pointer states survive (from Proposition [Einselection]).
  \item Dynamics reduces to classical mechanics (structural
  identification: classical trajectories are the $N \to \infty$ limit
  of RS ledger dynamics with instantaneous decoherence).
\end{enumerate}
\end{proposition}

\begin{remark}[Status of (3)]
Part~(1) is an immediate consequence of the decoherence time formula.
Part~(2) follows from the einselection argument.  Part~(3) is the
physical claim that the $N \to \infty$ limit of the quantum dynamics
produces Newtonian mechanics.  A rigorous derivation of this
limit---showing the RS Hamiltonian reduces to the Hamilton equations
of motion in the $N \to \infty$ limit---is an identified open problem.
\end{remark}

\begin{corollary}[Quantum--Classical Crossover]
A system is either quantum ($\tau_{\mathrm{dec}} > \tau_{\mathrm{obs}}$,
coherence observable) or classical ($\tau_{\mathrm{dec}} <
\tau_{\mathrm{obs}}$, coherence unobservable).  Setting
$\tau_{\mathrm{obs}} \sim 1$~s and $\tau_0 \sim 10^{-44}$~s, the
crossover occurs at:
\begin{equation}
  N_{\mathrm{cross}} \sim \frac{\ln(10^{44})}{\ln\varphi} \cdot \frac{1}{g}
  = \frac{44 \ln 10}{0.4812} \cdot \frac{1}{g}
  \approx \frac{210}{g}
\end{equation}
environmental modes.  For $g \sim 10^{-4}$: $N_{\mathrm{cross}}
\sim 2 \times 10^6$, corresponding to large molecules; for
$g \sim 10^{-2}$: $N_{\mathrm{cross}} \sim 2 \times 10^4$,
corresponding to quantum dots.
\end{corollary}

\section{Ehrenfest's Theorem}

Ehrenfest's theorem states that $\langle x \rangle$ and $\langle p
\rangle$ satisfy $d\langle x \rangle/dt = \langle p \rangle/m$ and
$d\langle p \rangle/dt = -\langle \nabla V \rangle$.  For general
potentials, $\langle \nabla V \rangle \neq \nabla V(\langle x
\rangle)$, so expectation values do not follow classical trajectories
exactly.  In RS, $J$-cost forces decoherence into pointer states
(position eigenstates) as $N \to \infty$.  The reduced density matrix
becomes diagonal in position; surviving components follow classical
trajectories.  Ehrenfest's theorem holds \emph{exactly} for the
effective classical dynamics because the system is a statistical
mixture of classical-like trajectories.  $J$-cost thus explains why
expectation values follow classical trajectories at macroscopic
scales.

\section{Mesoscopic Regime}

Between microscopic ($N \sim 1$--$10^3$) and macroscopic ($N \sim
10^{23}$) regimes lies the mesoscopic regime ($N \sim 10^4$--$10^8$).
Here $\tau_{\mathrm{dec}}$ is long enough for quantum coherence to be
observable.  RS predicts $N_{\mathrm{cross}} \approx 210/g$; for
$g \sim 10^{-5}$--$10^{-2}$, coherence persists from seconds down
to femtoseconds.  For $g \sim 10^{-4}$ (typical electromagnetic coupling
in mesoscopic systems), $N_{\mathrm{cross}} \approx 2 \times 10^6$.

\textbf{Fullerene interference.}  C$_{60}$ matter-wave
diffraction~\cite{Arndt} demonstrates coherence at $N \sim 10^3$
atoms.  RS places this near the lower mesoscopic boundary; weak
environmental coupling keeps $\tau_{\mathrm{dec}}$ observable.

\textbf{Quantum dots.}  Confined systems ($N \sim 10^4$--$10^6$)
exhibit quantized levels, tunneling, and coherent manipulation.
$J$-cost remains manageable; decoherence is tunable via size and
temperature.

\textbf{SQUID loops.}  Superconducting loops ($N \sim 10^{10}$--$10^{12}$
Cooper pairs) show coherent superposition of currents.  RS attributes
this to the collective state: the effective $N$ of environmentally
coupled degrees of freedom is reduced by the gap.

\section{Interpretive Comparison}

Table~\ref{tab:interpretations} compares how major interpretations of
quantum mechanics handle the quantum-to-classical transition and
decoherence.

\begin{table}[htbp]
\centering
\caption{Approaches to the quantum--classical border and decoherence.}
\label{tab:interpretations}
\begin{tabular}{@{}lp{4cm}p{4cm}@{}}
\toprule
Interpretation & Decoherence role & Quantum--classical border \\
\midrule
Copenhagen & \parbox[t]{4cm}{Collapse postulate; observer-dependent} & \parbox[t]{4cm}{Measurement cuts chain; no scale derivation} \\
Many-Worlds & \parbox[t]{4cm}{All branches exist; decoherence = branching} & \parbox[t]{4cm}{No objective border; appearance subjective} \\
Bohmian & \parbox[t]{4cm}{Pilot wave + particles; decoherence irrelevant} & \parbox[t]{4cm}{Particles always classical; wave guides deterministically} \\
GRW & \parbox[t]{4cm}{Spontaneous collapse; collapse causes decoherence} & \parbox[t]{4cm}{$\lambda \sim 10^{-16}$/s per particle; $N$-dependent amplification} \\
RS & \parbox[t]{4cm}{$J$-cost minimization; einselection from cost} & \parbox[t]{4cm}{$N \to \infty$ limit; $\tau_{\mathrm{dec}} \propto \varphi^{-Ng}$} \\
\bottomrule
\end{tabular}
\end{table}

RS uniquely derives decoherence timescale and pointer-state selection
from a single cost principle, without collapse postulates or new constants.

\section{Predictions}

RS predicts: (1) \textbf{Decoherence scaling:}
$\tau_{\mathrm{dec}} \propto \varphi^{-Ng}$; testable via rates vs.\
temperature or pressure.  (2) \textbf{Mesoscopic boundary:} For
$g \sim 10^{-6}$--$10^{-4}$, the crossover
$N_{\mathrm{cross}} \approx 210/g \sim 10^6$--$10^8$; virus
interferometry ($N \sim 10^7$--$10^8$) should show reduced
coherence.  Precision measurements of coherence visibility vs.\ $N$
would test this scaling law.  (3) \textbf{Pointer-state universality:}
Position emerges as pointer basis.  (4) \textbf{Minimal new parameters:}
Unlike GRW ($\lambda$, $r_C$), RS introduces only one phenomenological
coupling $g$ per environment type, expected to be derivable from the
RS Hamiltonian once the ledger--environment interaction is fully
specified.

\section{Conclusion}

The quantum-to-classical transition is smooth and automatic:
large $N$ forces rapid decoherence via quadratic $J$-cost scaling.
Classical mechanics emerges as the $N \to \infty$ limit.  No
additional postulates or interpretive choices are required.

\begin{thebibliography}{9}

\bibitem{Aczel1966}
J.~Acz\'el, \emph{Lectures on Functional Equations and Their Applications},
Academic Press, New York (1966).

\bibitem{RS_Axioms}
J.~Washburn, ``The Algebra of Reality: A Recognition Science Derivation
of Physical Law,'' \emph{Axioms} \textbf{15}(2), 90 (2026).
\texttt{doi:10.3390/axioms15020090}

\bibitem{Zurek}
W.~H.~Zurek, ``Decoherence, Einselection, and the Quantum
Origins of the Classical,'' Rev.\ Mod.\ Phys.\ \textbf{75},
715 (2003).

\bibitem{Arndt}
M.~Arndt \emph{et al.}, ``Wave--Particle Duality of
C$_{60}$ Molecules,'' Nature \textbf{401}, 680 (1999).
\end{thebibliography}

\end{document}
